\documentclass[../main.tex]{subfiles}
\begin{document}
%==========================================TIÊU ĐỀ TẬP========================================
\begin{table}[h]
  \begin{adjustwidth}{-1cm}{}
    \begin{tabular}{>{\centering\arraybackslash}p{6cm}|>{\raggedright\arraybackslash}p{12.5cm}}
      \multirow{5}{6cm}
      % Cột bên trái: ÉPISODE và số tập
      {\\[-40pt]
        \flaregothic\fontsize{20pt}{20pt}\selectfont \centering
        {\contour{errorthree}{\textcolor{black}{ÉPISODE}}}\color{black}\\
        \burbank\fontsize{72pt}{72pt}\selectfont
        \includegraphics[height=3.12359cm]{../../../../Image/Book?/number/num14.png}
        \\[18pt]
        % \flaregothic\fontsize{14pt}{14pt}\selectfont
        % \color{epattr}BIENVENUE, \uline{\color{Banpire} Mel} !
        % \flaregothic\fontsize{18pt}{18pt}\selectfont
        % \color{epattr}ÉPISODE PERDU
        \color{black}
      }
       &
      % Dòng thứ nhất: tên tập tiếng Nhật
      {\fotskip\fontsize{18pt}{18pt}\selectfont \ruby{何}{いつ}かの\ruby{記}{き}\ruby{憶}{おく}
      }                                    \\[6pt]
      % Dòng thứ hai: tên tập tiếng Anh
       & {\cafeteria\fontsize{18pt}{18pt}\selectfont Forgotten Memories}                                                \\[4pt]
      % Dòng thứ ba: tên tập tiếng Việt
       & {\vnmsans\fontsize{18pt}{18pt}\selectfont Ký ức bị lãng quên}                                                  \\[8pt]
      % Dòng thứ tư: Nhân vật xuất hiện
       & {\caviar{\fontsize{14pt}{14pt}\selectfont Track used: The Chaos (Guardians - Bob Bradley, Terry Devine-King)}}
      \\[6pt]
      % Dòng cuối cùng: Thời gian diễn ra của tập (thường sẽ là ngày phát hành)
       & {\continuum\fontsize{16pt}{16pt}\selectfont Ngày phát hành: 20/05/2022}                                        \\[8pt]
    \end{tabular}
  \end{adjustwidth}
\end{table}
%===========================================NỘI DUNG==========================================
\color{black}

\pov{Haragen Reiko}

Bước chân vội vã cuối cùng cũng đưa năm chúng tôi bắt kịp Akane-san - người đang chạy trước chúng tôi để lùng sục tìm Yuka-san - thì hoàn toàn mất hút. Cô ấy dừng lại nơi giao nhau giữa hai toa, tay vẫn siết chặt chiếc đèn pin nhỏ, ánh sáng loé lên loé xuống in bóng cô trên vách kim loại.

``May quá, cậu đến rồi!'' cô nàng tai thỏ thở phào, nắm chặt máy ảnh.

``Bọn mình lo muốn chết,'' cô nàng tóc nâu nói khẽ, mắt vẫn dán lên hành lang tối.

``Không sao,'' tôi đáp, kéo cao chiếc túi xách chéo vai của mình. ``Đi thôi nào.''

Không khí trong toa tàu vẫn dày đặc, nặng nề. Ánh sáng từ chiếc đèn pin mà Onii-sama cầm theo hắt thành một vệt mảnh, chiếu vào sàn đầy bụi và những hàng ghế lặng câm. Trên cao, đèn khẩn cấp nhấp nháy ánh đỏ yếu ớt, cứ vài giây lại chớp một lần, khiến cả hành lang dài hun hút như đang hít thở theo nhịp tim của con tàu.

Chúng tôi tiếp tục bước dọc hành lang, những âm thanh duy nhất vang lên là tiếng giày chạm xuống sàn kim loại và tiếng rít khe khẽ từ hệ thống điện đã tắt ngúm. Mỗi bước chân đều kéo theo một âm vang rỗng tuếch, như thể con tàu đang thở dài trong giấc ngủ đông.

Đi trước dẫn chúng tôi, Sora-san khẽ lẩm bẩm, giọng khàn như bị bụi bám: ``Con tàu có vẻ như dừng thật rồi\dots''

Không ai trong số chúng tôi đáp, nhưng tôi biết mọi người đều nghĩ như vậy. Cảm giác bất an dâng lên trong tôi, khiến tôi vô thức bước sát vào Onee-sama. Tôi thì thầm, đủ để chỉ Onii-sama và chị nghe thấy: ``Hy vọng đây không phải là điềm xấu\dots''

Sora-san thoáng ngoái lại, khóe môi nhợt nhạt, đáp như để trấn an cả chính mình: ``Mong là vậy\dots''

Bất ngờ, Uzuki-san giơ chiếc máy ảnh cầm tay của mình lên, ống kính loáng ánh sáng phản chiếu từ đèn pin. Gương mặt cô ấy không hề tỏ ra sợ hãi, mà ngược lại, có chút hứng thú. Vừa lia ống kính quanh toa, cô vừa cười khẽ: ``Trong này cũng có kha khá chỗ để làm vài kiểu ảnh đấy.''

\img{e2-6a}

Tôi nhăn mặt, khẽ thở dài: ``Mình ngạc nhiên là cậu vẫn còn nghĩ tới ảnh vào lúc này đấy\dots''

Cô ấy quay sang nhìn tôi, ánh mắt không hề bối rối:
``Cậu biết mà, Reiko-san. Mình từng gặp những thứ còn đáng sợ hơn nhiều. Thế nên, thành thật thì, chuyện này chưa đủ làm mình lung lay tinh thần đâu.''

Onee-sama liếc nhìn cô gái ấy, mày hơi nhíu lại: ``Những thứ còn đáng sợ hơn\dots\ ý cậu là gì?''

Cô mỉm cười, nụ cười mơ hồ trong ánh sáng chập chờn: ``Mình đã lang thang khắp thị trấn Aogami, cả ngày lẫn đêm, chỉ để chụp lại những gì gắn liền với nơi này. Có lần, giữa đêm, mình đi qua bờ sông và thấy một bóng đen lướt trên mặt nước. Rồi có khi, trong một trạm xăng bỏ hoang, mình nghe tiếng bước chân ngay sau lưng, dù rõ ràng chẳng có ai ở đó cả\dots''

Onii-sama gật gù, như thể đã nghe quen chuyện này nhiều lần, còn Onee-sama khẽ rùng mình, nắm chặt lấy tay tôi hơn.

Sora-san, ngược lại, tròn mắt nhìn Uzuki-san, giọng xen lẫn kinh ngạc lẫn khó tin: ``Cậu\dots\ thật sự dám đi một mình trong những chỗ như vậy sao?''

Cô nàng tai thỏ hạ máy ảnh xuống, vai nhún nhẹ: ``Nếu muốn có những bức ảnh đặc biệt, thì phải liều thôi. Với mình, chụp ảnh ở Aogami không đơn thuần là sở thích nữa, nó giống như một cách để ghi lại những gì đang dần biến mất theo thời gian.''

Không khí trong toa như chùng xuống sau câu nói ấy. Tôi thoáng liếc sang Sora-san, gương mặt cô ấy ánh lên sự ngạc nhiên lẫn thán phục, khác hẳn với sự dè chừng ban đầu. Với chúng tôi - những người đã biết Uzuki-san từ trước đó - thì chẳng lạ gì, nhưng thấy phản ứng của cô ấy, tôi chợt nhận ra: đối với người ngoài, niềm đam mê mãnh liệt này quả thật gây ấn tượng mạnh.

Tôi khẽ lên tiếng, giọng nhỏ nhưng đủ để cả nhóm nghe thấy: ``Cậu ấy luôn như vậy. Mỗi lần thấy thứ gì đang trôi vào quên lãng, cậu ấy lại như bị thôi miên.''

Onee-sama thở dài, nhẹ nhàng: ``Cũng nhờ vậy mà cô ấy mới giật giải thưởng đấy nhỉ?''

Onii-sama cười khẽ, đèn pin rung nhẹ theo nhịp tay: ``Và cũng nhờ vậy mà tôi có thêm vài cơn ác mộng.''

Uzuki-san phá lên cười, tiếng cười vang trong toa tàu như thể xua tan bóng tối: ``Nếu không có những bức ảnh ấy, làm sao mọi người thấy được thị trấn này đã thay đổi như thế nào, đúng không?''

Sora im lặng một lúc, rồi cất giọng nhỏ nhưng rõ ràng: ``Có lẽ tớ cũng muốn xem những bức ảnh đó, nếu được.''

Tôi mỉm cười, cảm thấy bầu không khí nhẹ nhõm hơn một chút. Dù bóng tối vẫn bao trùm, nhưng có điều gì đó trong ánh mắt Uzuki-san khiến tôi tin rằng: dù con tàu này dừng lại ở đâu, chúng tôi vẫn sẽ tìm thấy người bạn mà chúng tôi muốn tìm.

\ct

Ánh đèn đỏ khẩn cấp phía trên đầu vẫn chớp nháy, nhưng nhịp của nó dần trở nên bất thường, lúc sáng lúc tắt, như thể có bàn tay vô hình nào đang bóp nghẹt lấy nguồn điện. Tôi siết chặt tay Onee-sama, tim khẽ đập nhanh.

``Reiko, em ổn chứ?'' Onee-sama khẽ hỏi, giọng đầy lo lắng.

``Em\dots\ em không sao,'' tôi đáp, cố giữ bình tĩnh nhưng giọng vẫn run nhẹ.

Chúng tôi bước sang toa tiếp theo. Và ở đó, trong khoảng sáng đỏ mờ ảo, một bóng người ngồi bất động trên hàng ghế. Cô ấy cúi đầu, đôi mắt hé lộ ánh xanh đục như cá ươn. Đôi môi thì không ngừng mấp máy, lặp đi lặp lại những từ ngữ rời rạc.

\img{e2-6b}

Sora-san sững người, thốt lên: ``Kia chẳng phải là\dots\ Kaoru-san sao?''

``Cậu biết cô ấy sao?'' Onee-sama nghiêng đầu, đôi mắt mở to đầy ngạc nhiên.

``Ừ\dots\ Mình đã từng gặp cô ấy trước đây, cùng với Miku và Saya. Họ\dots\ giống như những người bạn cũ của mình. Bọn mình từng trò chuyện, từng đi lang thang quanh thị trấn với nhau. Kaoru-san khi ấy luôn tỏ ra quan tâm, như một chị cả trong nhóm.''

Cô ngừng lại, mắt khẽ chao đảo khi nhìn người đang ngồi trước mặt: ``Nhưng\dots\ hình ảnh đó bây giờ, thật sự không còn giống như xưa nữa.''

Uzuki lùi lại một bước, giọng xen chút ngập ngừng: ``Nhắc mới thấy, trông cô ấy có vẻ hơi\dots\ kỳ lạ. Tại sao cô ấy lại ở đây nhỉ\dots''

``Chúng ta nên cẩn thận,'' Onii-sama cảnh báo, tay siết chặt đèn pin.

Chúng tôi tiến lại gần hơn. Lúc này, tôi mới nghe rõ được tiếng lẩm bẩm đều đặn phát ra từ đôi môi tái nhợt ấy:

``\dots Đó cũng là tên của ngôi đền vùng này. Đó cũng là tên của ngôi đền vùng này\dots''

Một luồng lạnh chạy dọc sống lưng tôi. Tôi nhìn sang Onii-sama và Onee-sama, và nhận ra ánh mắt họ cũng đang tối lại. Quá quen thuộc. Quá giống với những gì chúng tôi đã gặp trước đây: những thực thể không còn là người, chỉ biết lặp lại một câu nói vô nghĩa.

Kaoru-san\dots\ không còn là Kaoru-san mà chúng tôi biết nữa.

Tôi nhớ lại lời cảnh báo của Shino-san: ``Nếu các cậu nhìn thấy những người từng gặp khi tới đây, tốt nhất đừng bắt chuyện. Họ không còn là người nữa.''

Uzuki-san ngập ngừng, dường như định lên tiếng hỏi: ``Kaoru-san\dots?''

``Đừng,'' tôi lập tức ngăn lại, lắc đầu, giọng gấp gáp nhưng dứt khoát. ``Đi tiếp thôi.''

``Phải đấy, Uzuki-san. Tin mình đi,'' Sora-san cũng lên tiếng, giọng trầm hơn thường ngày.

Cô nàng tai thỏ cắn môi, đôi mắt thoáng hiện vẻ khó hiểu và lo lắng, nhưng cuối cùng cũng gật đầu, ôm chặt chiếc máy ảnh rồi lặng lẽ bước theo chúng tôi, mặc kệ cho Kaoru-san - đúng hơn là cái xác của cô ấy - lặp đi lặp lại câu nói như thể đang độc thoại một mình.

Khi ba anh chị em tôi đi ngang qua, tôi cảm nhận được rõ rệt đôi mắt vô hồn ấy chuyển động, lia theo từng bước chân của chúng tôi. Trong khoảnh khắc, hơi thở tôi bỗng chốc cảm thấy nặng nề. Rồi, trước mắt tôi, thân hình Kaoru-san dần tan ra thành một bóng khói đen, vặn vẹo như thể đang tự nuốt chửng chính mình, thứ hình dạng quá quen thuộc với chúng tôi từ những vòng lặp trước.

``Cô ấy\dots\ biến mất rồi,'' Onee-sama thì thầm, giọng đầy nỗi kinh hoàng.

Mồ hôi lạnh túa ra trên trán. Chúng tôi đã thấy quá nhiều lần cảnh tượng này, vậy mà mỗi lần vẫn không tránh khỏi rùng mình.

``Chúng ta phải tiếp tục thôi,'' Onii-sama nói, giọng đầy quyết tâm nhưng tay vẫn run nhẹ.

Chúng tôi bước đi, nhưng tôi không dám ngoái lại. Mỗi bước chân xa dần, tôi vẫn cảm thấy một luồng hơi lạnh níu kéo sau lưng, như thể cái bóng đen vừa tan ấy vẫn còn nguyên hình hài, lặng lẽ bám theo chúng tôi trong bóng tối.

``Cảm giác này\dots\ giống hệt lần Miku-san hiện trên cửa sổ hồi nãy, mắt trống rỗng, miệng mấp máy những âm thanh không thành lời,'' tôi nói nhỏ, giọng đầy nỗi sợ hãi.

Tôi siết chặt tay Onee-sama hơn, nhưng không dám nói. Chỉ có tôi biết, cái bóng ấy chưa buông tha.

\ct

Chúng tôi vừa rời khỏi toa có Kaoru, thì ánh sáng lờ mờ phía trước hé lộ một dáng người quen thuộc. Akane-san. Cô vẫn bước đi chậm rãi nhưng cẩn trọng, mắt dán vào khoảng tối phía trước, dường như chẳng mảy may để ý đến bầu không khí rợn ngợp quanh mình.

``A, Akane-san!'' Uzuki nhận ra đầu tiên, giọng bật lên khe khẽ.

Cô vội tiến lên vài bước, giơ tay như muốn chạm vào vai đối phương, rồi tiếp tục: ``Cái chuyện vừa rồi ấy--''

Cô nàng tai chó bất ngờ quay lại. Khuôn mặt cô không tức giận, nhưng ánh mắt và biểu cảm cứng lại như thép. Giọng nói trầm xuống, lạnh băng, đủ khiến da thịt tôi nổi gai: ``\dots Gì nữa? Tôi đang đi tìm Yuka đấy. Đừng làm phiền tôi.''

\img{e2-6c}

``Nhưng mà\dots'' Uzuki khựng lại, đôi môi mấp máy như muốn nói thêm, nhưng cô ấy đã ngoảnh mặt đi, tiếp tục sải bước, mặc cho lời cô còn bỏ lửng.

``Cô ấy\dots\ thay đổi rồi.''

Tôi thoáng thấy nét buồn trên gương mặt Sora-san, còn bản thân thì đưa tay chạm nhẹ vai cô tai thỏ, lắc đầu nhẹ thay cho lời nhắc.

``Trước mắt thì chúng ta nên giúp Akane-san tìm thấy Yuka-san trước đã,'' cô nàng tóc nâu dài lên tiếng, giọng nhẹ nhưng dứt khoát.

``Ừ, chúng ta không thể để cô ấy một mình,'' tôi đồng tình, nhìn theo bóng cô nàng tai chó khuất dần trong bóng tối.

Tất cả đều gật đầu. Không ai muốn gây thêm xung đột trong tình thế này. Thế là cả nhóm lặng lẽ bước theo sau Akane-san, vừa đưa mắt quan sát xung quanh, vừa dè chừng không dám làm động lớn.

``Mọi người\dots\ có cảm thấy gì lạ không?'' Uzuki-san thì thầm, giọng đầy lo âu.

``Cứ như\dots\ có thứ gì đang theo dõi chúng ta vậy,'' Sora-san đáp, mắt không ngừng liếc nhìn xung quanh.

Bỗng nhiên\dots

*CỘP CỘP CỘP*

``\textbf{Gyaaaah!!!}'' tất cả sáu người chúng tôi hét lên.

\img{e2-6d}

Tiếng lục cục mơ hồ vang lên từ ngoài thân tàu. Một nhịp, rồi hai, rồi ba. Mỗi tiếng như gõ thẳng vào lồng ngực.

``Tiếng gì vậy!?'' tôi bật thốt, giọng run run.

``Nghe như có ai đó\dots\ gõ ở bên ngoài,'' Onii-sama đáp, từng chữ nặng nề.

``Không thể nào\dots\ chúng ta đang trên tàu mà,'' Onee-sama nói, nhưng giọng cũng đầy nghi ngờ.

``Có thể là gió thôi?'' Uzuki thì thầm, nhưng ngay cả cô ấy cũng không tin nổi lời mình nói.

``Không, gió không tạo ra nhịp đều đặn như vậy,'' Akane lắc đầu, mắt vẫn dán vào cửa kính.

Chúng tôi chầm chậm đưa mắt về phía ô cửa kính. Và ở đó\dots

Một cô gái hiện ra, mặt áp sát lên khung cửa, mái tóc vàng ngắn rối bời, cài chiếc kẹp tóc hình con dơi. Bàn tay cô liên tiếp đập vào cửa kính, tạo nên chuỗi âm thanh lục cục dồn dập. Xung quanh, dấu bàn tay loang loáng xuất hiện chằng chịt, như thể đã từng có hàng chục, hàng trăm lần cô lặp đi lặp lại hành động đó.

``Nè nè\~{}\dots\'' cô cất tiếng, chậm rãi, nhẹ nhàng đúng với giọng điệu mà chúng tôi từng biết. Nhưng giờ đây, âm sắc ấy vỡ ra thành sự lạnh lẽo.

``Ta cùng chơi thôi nào\~{}''

\img{e2-6e}

``Saya-chan?'' Uzuki-san là người buột miệng đầu tiên, giọng run run. ``Cô ấy\dots\ sao lại ở ngoài cửa kính chứ?''

``Tại sao\dots\ cậu ấy lại hành xử như vậy?'' Akane-san cau mày, nghiêng đầu, giọng cô lạc đi vì sốc.
``Không giống Saya-chan chút nào\dots''

``Cô ấy\dots\ khả năng cao không còn là Saya-san nữa,'' tôi thì thầm, cố giữ bình tĩnh. ``Giống như Kaoru-san vừa rồi vậy.''

``Reiko-san?'' cô nàng tai thỏ quay sang nhìn tôi, đôi mắt mở to đầy hoảng loạn. ``Ý cậu là sao chứ?''

``Chúng ta phải rời khỏi đây,'' Onii-sama nói, tay siết chặt đèn pin. ``Không được nhìn vào mắt cô ấy.''

``Nhưng cô ấy đang ở ngoài kia mà!'' cô nàng tai thỏ phản đối, giọng gần như khóc. ``Chúng ta không thể bỏ mặc Saya-san được!''

``Cậu ấy đã không còn tồn tại nữa rồi,'' Onee-sama thì thầm, giọng đầy đau đớn. ``Chỉ còn lại vỏ bọc đang lặp lại ký ức thôi.''

Ánh đèn đỏ khẽ nhấp nháy, in bóng năm chúng tôi dài trên sàn kim loại. Không ai bước thêm bước nào, như thể mặt kính vừa trở thành bức tường vô hình ngăn cách hai thế giới. Gió lạnh từ khe cửa rít lên, cuốn theo mùi sắt gỉ và hơi ẩm của đêm khuya, khiến tôi tự nhiên nghẹn thở. Trong khoảnh khắc ấy, tiếng tim tôi đập mạnh đến mức nghe như trống ngực ai đó đang gõ thình thịch bên tai.

``Ta cùng chơi thôi nào\~{}'' Saya vẫn nở nụ cười ma mị ấy, đôi bàn tay trắng bệch dập mạnh vào ô cửa kính, nhịp gõ vang vọng khắp khoang tàu như tiếng kim loại vỡ vụn trong đầu.

``Mọi người ơi\dots\ tại sao không mở cửa cho mình vậy\~{}?''

``Im lặng,'' tôi rít lên, kéo Uzuki-san lùi lại. ``Đừng trả lời.''

Thế rồi, từ bên ngoài, một làn khói đỏ cuồn cuộn kéo tới, ôm trọn lấy thân hình gầy guộc của cô ấy. Tiếng Saya như tan ra thành tiếng gió rít.

Cả thân người Saya bị khói đỏ bao phủ, rồi từng chút một tan thành bụi li ti. Những hạt bụi đỏ rực vỡ nát trong ánh đèn khẩn cấp nhấp nháy, trước khi bị một cơn gió lạnh lùa qua thổi tan biến. Khi lớp khói lắng xuống, cửa sổ chỉ còn trơ trọi bóng tối bên ngoài. Không còn Saya nữa. Tựa như cô chưa từng tồn tại.

\gif{e2-6f}{1}{116}

``Như vậy là chắc chắn rồi,'' anh tôi khẽ gằn giọng, rất thấp.

``Em cũng nghĩ thế. Cậu ấy đã không còn là Saya mà chúng ta biết nữa,'' chị tôi gật khẽ, đôi mắt lặng thinh.

``Shino-san đã đúng\dots'' tôi chỉ biết siết chặt hai tay, thì thầm.

Không ai trong chúng tôi dám nói điều này với Akane-san và Uzuki-san. Chúng tôi không dám chắc họ sẽ nghĩ gì, và liệu họ có tin lời chúng tôi nói không.

Vậy nên, tốt nhất là chúng tôi vẫn nên giữ im lặng.

``Đ-Đó có phải là\dots\ Saya-san không vậy?'' cô nàng tai chó vẫn còn choáng váng, lắp bắp.

``Mình nghĩ\dots\ là như vậy,'' cô nàng tai thỏ mặt tái đi, nhưng vẫn gượng đáp, giọng chẳng hề chắc chắn.

``Chúng ta phải tiếp tục thôi,'' Sora-san lên tiếng, giọng khàn đặc.

Trong khoảnh khắc ấy, tiếng còi tàu vang lên chói tai, xuyên qua sự im lặng chết người. Toa tàu rung lắc nhẹ, rồi chuyển động trở lại. Ánh đèn đỏ trên trần nhấp nháy liên hồi, hắt xuống sàn những vệt sáng và bóng chồng chéo khiến tôi chỉ càng thấy như đang bị mắc kẹt trong cơn ác mộng.

``Mình\dots\ mình phải tìm Yuka\dots'' Akane-san lúc này mới giật mình như bừng tỉnh khỏi cơn mê. Cô bặm môi, thì thầm gần như tự nói với mình.

Không đợi thêm một giây, cô lao nhanh về phía toa kế cận, để mặc chúng tôi đứng sững lại sau lưng.

``Cô ấy sẽ ổn thôi,'' tôi nói, nhưng ngay cả tôi cũng không tin nổi lời mình.

``Chúng ta nên đuổi theo,'' Onii-sama đề nghị, mắt vẫn dán vào bóng tối nơi cô ấy vừa biến mất.

``Uzuki-san\dots' Sora-san bây giừo mới dám thở phào, đôi vai khẽ run lên. Cô nhìn sang cô bạn tai thỏ, thấy cô đang cúi xuống nhặt lại chiếc máy ảnh rơi từ trước, đôi bàn tay còn run rẩy.

``Trông cậu có vẻ mệt đấy. Cậu ổn chứ?'' cô nhẹ giọng, ánh mắt lo lắng.

``Ừ\dots\ Hơi khó thở một chút, nhưng mà mình vẫn ổn thôi,'' Uzuki-san gượng cười nhạt.

``Cậu có muốn một chút nước không? Uống đi, sẽ thấy khá hơn đấy,'' tôi lục chiếc túi xách, lấy chai nước đưa ra.

``Cảm ơn nhé.''

Cô nhận lấy, uống một hơi dài, rồi khẽ ngửa mặt thở ra, như cố gắng trấn tĩnh lại sau những sự kiện vừa rồi.

``Chúng ta không thể dừng lại được,'' Onee-sama nhắc nhở, mắt nhìn về phía trước.

``Đúng vậy, cứ đứng mãi ở đây cũng không giúp ích gì,'' tôi đồng tình.

Onii-sama ra hiệu bằng ánh mắt, Onee-sama cũng gật đồng tình. Thế là cả nhóm không nói thêm gì nữa, chỉ lẳng lặng bước tiếp, nhanh chóng đuổi theo Akane về phía toa tàu kế cận.

Trong bóng tối ngập ngụa ánh đỏ, từng tiếng bước chân nện xuống sàn vang vọng, như tiếng trống dồn dập báo hiệu cho một điều chẳng lành sắp sửa xảy đến.

\ct

Chúng tôi bước đi trong im lặng, bước chân vang lên những tiếng kim loại lách cách nhỏ xíu trong khoảng lặng nặng nề. Ánh đèn đỏ nhấp nháy hắt lên vách toa, khiến bóng chúng tôi dập dờn như những bóng ma bị kẹt trong đường hầm vô tận. Không khí lạnh lẽo quấn quanh cổ họng, mỗi hơi thở đều nặng trĩu như thể đang hút đầy bụi mù tanh của thời gian bỏ hoang.

Cho đến khi Onee-sama đột nhiên khựng lại, mắt dán về phía một chỗ ngồi bên cạnh cửa sổ. Cô chậm rãi nâng tay, chỉ vào khoảng không mờ ảo dưới ánh đèn đỏ, giọng thì thầm chỉ đủ chúng tôi nghe: ``Ơ\dots\ Quyển sách kia là\dots'' Giọng chị xen lẫn chút cảnh giác, như thể vừa phát hiện ra một bẫy ngầm.

Tôi cùng mọi người nhìn theo. Quả thật, có một quyển sách dày cộp đang mở hé, nằm vất vưởng như thể có ai vừa đặt xuống rồi bỏ đi vội. Bìa cứng màu xám tro đã bong tróc vài chỗ, lộ lớp bìa lót màu nâu sẫm như máu khô. Ánh đèn đỏ nhấp nháy hắt lên bìa sách, khiến nó càng trở nên khác thường, như một con mắt đang chớp chớp nhìn xuyên qua linh hồn.

Sora-san nuốt nước bọt, khẽ nghiêng đầu, nói với giọng thì thào và khô khốc: ``Có ai đó\dots\ đã đọc nó gần đây lắm.''

Chúng tôi tiến lại gần, bước chân nhẹ nhàng. Cô đưa tay cầm nó lên, bụi mỏng phủ ra bay tản mác, lượn lờ dưới ánh đèn như những con sâu bụi bẩn đang bỏ chạy. Cô nhẹ nhàng lật nhìn, ngón tay run run chạm vào mép giấy đã ngả vàng.

Trên hai trang đang giở sẵn:

Trang bên trái in đậm hàng chữ: ``Kỷ yếu tốt nghiệp trường Trung học Aogami, niên khóa 2015-2018'', bên dưới là biểu trưng của trường cùng một bức ảnh chụp ngôi trường quen thuộc.

Trang bên phải là lời cảm ơn và chữ ký của hiệu trưởng, nắn nót từng nét, nhưng mực đã loang một vệt dài, như thể có người từng khóc lên trang giấy.

\img{e2-6g1}

Sora-san khẽ thốt lên, giọng nhỏ đến mức gần như tan trong không khí: ``\dots Đây là\dots''

Uzuki ghé sát, nghiêng đầu, mắt không rời: ``Cuốn\dots\ kỷ yếu ư?''

Onii-sama nhíu mày, gằn giọng thấp, như thể sợ đánh thức ai đó: ``Tại sao nó lại ở đây nhỉ? Một thứ như thế mà nằm trong toa tàu bỏ hoang này\dots\ Không hợp lý chút nào.'' Anh khẽ đưa tay vuốt mép sách, lòng bàn tay dính một lớp bụi mịn như tro.

Sora-san im lặng, hít một hơi thật sâu, rồi run tay lật sang trang kế bên. Tiếng giấy ken két vang lên như tiếng xương khô gãy.

Một danh sách học sinh tốt nghiệp hiện ra. Nhưng\dots\ tôi thấy rõ, phần lớn các khuôn mặt đều bị bôi đen xóa nhòa, không chỉ bằng mực, mà còn bằng những vệt xước như móng tay cào lên ảnh. Cả tên dưới ảnh cũng không còn đọc được, chỉ còn lại những khoảng trắng đen sì như hố sâu.

Chỉ có sáu gương mặt là hiện rõ: sáu khuôn mặt tươi cười, nhưng nụ cười họ trông có vẻ khá tự nhiên, càng làm cho chúng tôi rùng mình.

\img{e2-6g2}

Người cầm quyển lùi lại một chút, mắt mở to, đồng tử co lại: ``Hể\dots\ Cái gì thế này\dots?''

Tôi không kìm được mà thì thào, giọng khô khốc: ``Hầu như tất cả mọi người đều bị tô đen mặt luôn kìa\dots\ Cả tên của họ cũng vậy. Như thể\dots\ họ chưa từng tồn tại.''

Onii-sama và Onee-sama liếc nhìn kỹ, rồi lần lượt chỉ vào từng người, giọng họ trở nên trầm đục, như thể đang đọc bia mộ:

``Đây\dots\ cô gái tóc vàng, đeo nơ xanh\dots\ ghi là Shiratori Hotaru,'' anh nói nói, ngón tay dừng lại trên đôi mắt cô gái.

``Còn cô tóc tím, khuyên tai, tóc hai bím xoắn ốc\dots\ tên là Tokiwa Kana.'' Onee-sama tiếp lời, giọng mỗi lúc một nhỏ.

``Cô gái tóc trắng, cài hào quang thiên sứ\dots\ Amano Kanade.''

``Người tóc hồng, mắt hai màu xanh lá mạ và hồng\dots\ Himekawa Mitsuki.''

``Cô tóc đen, dần xuống đỏ\dots\ Kaizuka Yuki.''

Và rồi cả ba chúng tôi cùng dừng lại ở bức ảnh cuối cùng, khuôn mặt nằm ở góc dưới cùng, như thể bị ép vào khung ảnh bằng một lực vô hình.

Một cô gái tóc nâu dài, đôi mắt xanh sắc lạnh, chiếc cài tóc đỏ lệch về một bên. khuôn mặt trông quá quen thuộc với chúng tôi, quen đến mức khiến tim tôi nhói lên. Dưới đó, hàng chữ ghi rõ:

``Misora Shino.''

\img{e2-6h}

Khoảnh khắc cái tên ấy được đọc thành tiếng, Sora-san bỗng chốc khựng lại. Cô đứng im như tượng đá. Đôi mắt cô trống rỗng trong thoáng chốc, rồi khẽ lẩm bẩm, giọng như từ hầm mộ vọng lên: ``\dots Kana-san\dots\ Kanade-san\dots'' Cô nói đi nói lại hai cái tên ấy, như thể đang cố gọi họ về từ cõi chết.

Uzuki-san giật mình, nắm chặt tay cô ấy: ``Hai cái tên trong cuốn kỷ yếu đó sao? Cậu biết họ à? Cậu\dots\ nhớ ra gì rồi?''

\pov{Sora}

Tôi ôm lấy thái dương, sắc mặt trắng bệch, môi run bần bật: ``Tớ\dots\ Tớ không biết\dots\ Nhưng\dots\ đầu tớ\dots\ như bị\dots\ búa đập\dots'' tôi thở gấp, từng luồng khí lạnh phả ra từ miệng.

Reiko-san thoáng khựng lại. Trong khoảnh khắc, tôi như bị hút vào dòng ký ức nào đó, nhưng cơn đau trong đầu không cho phép tôi quan tâm thêm.

Một thoáng, trong đầu tôi chợt hiện ra hai gương mặt lạ mà quen: một cô gái tóc tím, buộc hai bím xoắn tinh nghịch, và một cô gái mái trắng-xám, có vệt xanh dương sáng nổi bật, trên tóc cài phụ kiện hình vầng quang.

Ký ức mờ nhòe bỗng ùa về. Tôi thấy mình ngày đó, ngồi bên khung cửa sổ, đôi mắt lặng lẽ nhìn xuống sân trường, lòng nặng trĩu bởi mặc cảm, một học sinh mới, lạc lõng giữa ngôi trường Aogami xa lạ. Rồi tiếng gọi trong trẻo vang lên:

``Này, này, bạn học sinh mới chuyển đến ơi!''

``Cậu đã bắt đầu quen với ngôi trường Aogami này rồi chứ?''

Tôi giật mình ngẩng lên. Là họ\dots\ là họ! Kana và Kanade. Nhưng, tại sao? Tại sao hình ảnh ấy lại ở đây, ngay lúc này?

\img{e2-6j}

Một nhói buốt như lưỡi búa bổ thẳng xuống đầu, nén chặt óc tôi. Cơn đau dữ dội khiến tôi nghẹt thở, như có hàng ngàn mũi kim cắm vào thái dương. Tôi muốn hét lên, muốn xua đi hình ảnh ấy, nhưng chúng càng khắc sâu thêm.

Tại sao\dots\ tại sao họ lại ở đây?

\pov{Haragen Reiko}

Ngay sau câu đó, cô khụy gối, hai tay ôm chặt đầu, móng tay cào vào da thịt. Một cơn đau dữ dội dường như xé toạc nhận thức của cô, khiến Uzuki hoảng hốt cúi xuống đỡ lấy, giọng thét lên: ``\textbf{Sora-chan!? Cậu sao vậy!? Đừng làm vậy!}''

Tôi và hai anh chị cũng nhào tới. Onee-sama đặt tay lên vai cô ấy, giọng dịu đi như ru trẻ nhỏ:
``Bình tĩnh nào\dots\ hít thở chậm thôi\dots\ Một\dots\ hai\dots\ từ từ thôi\dots''

Tôi thì thầm, cố giữ bình tĩnh cho Uzuki-san, giọng run run: ``Không sao đâu\dots\ Sora-san sẽ qua được thôi. Cô ấy\dots\ mạnh mẽ hơn chúng ta tưởng.''

Onii-sama gật gù, ánh mắt nghiêm trọng nhưng đầy thấu hiểu. Đây chính là điều mà Shino-san từng nhắc đến\dots\ nhưng dường như chưa đủ mạnh để phá vỡ cô ấy hoàn toàn.

\dots

\dots

Một lúc sau, cô ấy run rẩy ngẩng đầu, trán lấm tấm mồ hôi lạnh, giọng yếu ớt như sợi chỉ tơ: ``Tớ\dots\ Tớ ổn. Chỉ là\dots\ nhói một chút thôi.''

Uzuki-san thở phào, nhưng ánh mắt vẫn đầy lo lắng, tay không dám rời khỏi vai bạn: ``Cậu có sao không? Có cần nghỉ không?''

``Thật đấy, tớ ổn rồi\dots'' Sora mấp máy, cố nở một nụ cười gượng, nhưng nụ cười ấy méo xệch, như thể sắp vỡ.

Cô bạn tai thỏ mím môi, rồi siết chặt tay mình, lấy lại tinh thần. Cô quay sang chúng tôi, giọng trở nên cứng cỏi: ``Dù gì thì\dots\ chúng ta phải đi theo Akane-san thôi. Không thể để cô ấy một mình được. Cô ấy\dots\ cũng đang vỡ vụn mất rồi.''

\img{e2-6i}

Nói rồi, cô dìu Sora-san đứng dậy, từng bước đi chậm rãi như đang bước trên mặt hồ thủy tinh vỡ.

``Uzuki-san,'' tôi quay lại, giọng nhỏ nhưng đủ nghe, ``cậu đi trước đi. Chúng mình sẽ bắt kịp sau. Để bọn mình thở một lát.''

``Ừm, mình biết rồi,'' cô bạn tai thỏ đáp lại.

Cả hai nhanh chóng đi về phía trước, bóng họ loang dần trong ánh đèn đỏ nhấp nháy.

Tôi, Onii-sama và Onee-sama vẫn đứng nán lại. Trong im lặng, cả ba trao đổi ánh nhìn, đều hiểu rằng chuyện vừa xảy ra không hề đơn giản. Cuốn kỷ yếu vẫn mở hé trong tay tôi, trang giấy cuối cùng phập phồng như đang thở. Một mùi khét lạ thoang thoảng, mùi của ký ức bị thiêu cháy.

``Nhưng mà\dots'' tôi chợt lên tiếng, ngón tay vuốt nhẹ mép giấy như sợ nó tan biến, ``em có cảm giác\dots\ cuốn này đang kêu gọi chúng ta.''

Onii-sama nheo mắt, đèn pin hắt xuống lớp bìa mờ như phủ tro: ``Kêu gọi kiểu gì?''

``Như thể\dots'' tôi nuốt khan, ``như nó muốn được đọc tiếp. Mỗi lần em định khép lại, dường như có một thế lực nào đó ép nó phải mở ra ạ.''

Onee-sama rùng mình, đặt tay lên mu bàn tay tôi: ``Vậy thì\dots ta cùng đọc. Nhưng lần này, cả ba cùng lật.''

Anh gật, giọng khàn: ``Nếu đây là cái bẫy, thì cả ba người chia lưỡi dao cũng đỡ đau hơn.''

Ba cái bóng dồn lại, ba hơi thở gắn vào nhịp tim của trang giấy. Tôi kéo nhẹ, tiếng giấy rít lên như màn cửa cũ trong rạp chiếu bóng ma. Một luồng gió lạ tỏa ra từ khe giữa các trang, mang theo mùi mực in lâu năm và mùi máu khô thoang thoảng.

``Không chỉ là ký ức bị thiêu,'' chị khẽ nói, mũi khẽ động, ``còn có mùi\dots\ kim loại nữa.''

``Mùi của tên tuổi bị rút ra khỏi xương,'' người anh thì thầm, mắt không rời khỏi trang sắp lật, ``hoặc của ai đó đang cố nhớ lại mình từng là ai.''

Tôi siết tay hai anh chị: ``Em sợ\dots\ nhưng em cũng tòa. Giống như có dây thừng vô hình buộc vào cổ tay, kéo chúng ta xuống hầm ký ức.''

Onee-sama mỉm cười buồn: ``Vậy thì để anh chị đi trước. Nếu hầm đó có quỷ, nó phải bước qua xác chúng ta mới tới em.''

Onii-sama hít một hơi thật sâu, đặt ngón tay lên mép trang kế tiếp: ``Ba\dots\ hai\dots''

``Một,'' tôi thì thầm, và cả ba cùng lật.

Tôi run run lật sang trang tiếp theo. Tim tôi như thắt lại khi nhìn thấy lớp khác được in trong cuốn kỷ yếu. Vẫn cùng một kiểu trình bày: phần lớn khuôn mặt bị bôi đen xóa mờ, chỉ để lại vỏn vẹn ba gương mặt sáng rõ.

Và rồi\dots\ tôi chết lặng.

``Đ-Đừng nói với anh là\dots'' Onii-sama buột miệng, bàn tay anh siết chặt cuốn sách đến mức run lên.

Tấm ảnh đầu tiên, rõ ràng là anh - Haragen Kyuen - nhưng bên dưới lại ghi: ``Hikawa Kuon.''

Tấm thứ hai, gương mặt chẳng thể nhầm vào đâu của Haragen Saikyou, tên lại là: ``Hikawa Kaigou.''

Tấm cuối cùng, chính là tôi, nụ cười hơi gượng, mái tóc đen với hai bím thấp, đôi mắt xanh-đỏ quen thuộc. Nhưng dòng chữ bên dưới không phải ``Haragen Reiko'' mà là ``Hikawa Suzuki.''

Tôi lập tức đưa tay che miệng, cảm giác choáng váng ùa đến. ``Sao\dots\ sao lại là tên này\dots? Em\dots\ em không hiểu\dots''

Chị tôi giật lấy cuốn kỷ yếu từ tay anh, lật qua lật lại, mặt tái đi. ``Cái này không phải trò đùa. Mặt chúng ta là thật\dots\ nhưng họ ghi họ là Hikawa. Lại còn Kuon, Kaigou, Suzuki nữa\dots''

Anh tôi thì lẩm bẩm, giọng nghẹn lại như bị choáng: ``Nếu\dots\ nếu đây mới là tên thật\dots\ thì tất cả ký ức của chúng ta, gia đình Haragen của chúng ta\dots\ đều là giả sao?''

Trong đầu tôi bùng lên một mớ hỗn loạn. Từng ký ức hiện ra: tôi nhớ những bữa ăn gia đình, tiếng cười của Onii-sama khi trêu chọc em mình, vòng tay của cha mẹ khi tôi còn bé. Những điều đó\dots\ đều quá chân thật. Không thể nào chỉ là lớp vỏ dối trá. Nhưng cuốn kỷ yếu này lại như một chứng cứ tàn nhẫn, lạnh lùng chỉ ra rằng chúng tôi đã sống với một cái tên không thuộc về mình.

Tôi thấy mình như bị kéo tuột khỏi mặt đất, mất đi điểm tựa duy nhất. Nếu cái tên có thể thay đổi, liệu ký ức của tôi có thật không? Liệu tình cảm gia đình giữa chúng tôi có phải cũng chỉ là một ``kịch bản'' được sắp đặt?

Tôi không kìm được nữa, nước mắt cứ thế trào ra. Tôi lao tới ôm chặt lấy anh chị. ``Không\dots\ không thể nào\dots\ Anh chị là gia đình của em. Nếu tên khác đi\dots\ thì em phải tin vào cái gì nữa chứ?''

Khoảnh khắc đó, ký ức chập chờn lướt qua đầu tôi. Cái đêm trước vòng lặp thứ ba, khi Shino-san đã từng nhắc đến cái tên ``Hikawa\dots'' đó, ngay khi chúng tôi bước qua cánh cổng.

Lúc đó tôi không hiểu, chỉ thấy là một câu nói bí ẩn. Nhưng giờ đây, khi cái tên xuất hiện ngay trước mắt, tôi rùng mình.

Onee-sama cắn môi, giọng trầm xuống: ``Vậy ra\dots\ Shino-san đã biết từ trước. Cậu ấy cố tình nhắc đến cái tên Hikawa đó.''

Tôi siết chặt tay hai anh chị, toàn thân run lẩy bẩy nhưng vẫn cố gắng cất tiếng: ``Dù\dots\ dù thế nào đi nữa, em không tin ký ức của chúng ta chỉ là giả. Em không muốn tin. Em chỉ biết chắc một điều\dots\ chúng ta phải tìm ra sự thật. Phải hỏi thẳng Shino-san, hoặc bất kỳ ai biết điều này.''

Chị nhìn tôi, mắt rực sáng như bùng lên quyết tâm. ``Reiko nói đúng. Nếu có bí mật đằng sau cái tên này, thì chúng ta càng phải đi đến tận cùng.''

Anh tôi siết nắm tay, hít một hơi thật sâu, như trấn áp cơn run rẩy trong người. ``Ừ\dots\ dù là Kuon, Kaigou hay Suzuki\dots\ chúng ta vẫn là anh em. Không ai có thể phủ nhận điều đó. Hãy đi tìm câu trả lời.''

Tôi lau nước mắt, cố gắng nở một nụ cười yếu ớt. ``Vâng\dots\ Chúng ta sẽ làm sáng tỏ tất cả.''

Không khí trong toa tàu chợt nặng nề khác lạ, như đang ép chặt lấy lồng ngực. Nhưng giữa sự hoang mang, một sợi dây kết nối vô hình lại càng siết chặt chúng tôi, ba anh chị em Haragen\dots\ hay là Hikawa? Tôi không dám chắc.

Đúng lúc ấy, một tiếng cạch vang khẽ. Từ cuốn kỷ yếu, một vật mỏng, cứng rơi xuống sàn, lăn tròn rồi dừng lại ngay trước chân tôi. Tôi cúi xuống nhặt lên, một khung kim loại nhỏ, có hình chữ nhật rỗng giữa, những cạnh khít đến mức như được chế tạo để ôm lấy một linh kiện khác.

``Đây có lẽ là\dots\ công-xon?'' Onii-sama thì thào, ánh mắt thoáng ngạc nhiên khi một vật như vậy đột ngột xuất hiện.

``Có vẻ là như vậy,'' Onee-sama nắm chặt tay, giọng run nhẹ.

Tôi khựng lại: nó trông chẳng khác gì một cái khớp nối, một bộ phận nằm giữa mặt trước và mặt sau đồng hồ. Giống như một ``công-xon,'' giữ cho tất cả các mảnh rời rạc liên kết thành một khối hoàn chỉnh. Chúng tôi đã có pin, bảng mạch, lớp vỏ\dots\ và giờ đây, mảnh ghép này khiến chiếc đồng hồ tưởng như sắp sửa hiện hữu trọn vẹn trong tay.

``Mang hết ra đây,'' anh ra lệnh, quỳ gối xuống sàn. Từ túi áo khoác và túi quần, anh rút hai dây đeo đen nhánh còn vương bụi. ``Dây của anh.''

Chị tôi đặt hai nửa vỏ kim loại xuống cạnh, mặt trong khắc dòng chữ ``Gear Frontier'' vẫn còn nguyên vẹn. ``Vỏ của em nữa.''

Tôi lần lượt đặt cục pin mỏng, tấm màn hình nứt và bảng mạch chủ. Ba món của ba người, giờ nằm sát nhau như đang chờ giây phút tái sinh.

``Ta cùng ghép thử nào,'' tôi thì thầm. Cả hai anh chị đều cùng nhau gật đầu.

Onii-sama gắn bảng mạch vào nửa vỏ sau, tay cậu run đến mức vít nhỏ rơi lách cách. ``Khớp rồi\dots''

Onee-sama ấn màn hình xuống, mặt kính cũ kỹ phản chiếu gương mặt tái nhợt của chúng tôi. ``Vẫn còn khe hở\dots''

Tôi đặt cục pin vào khe, cảm giác lạnh toát từ kim loại khiến tôi rùng mình. ``Giờ thì công-xon.''

Anh cầm mảnh khớp nối bằng hai ngón tay, hít một hơi thật sâu rồi trượt nó vào khe hở giữa hai nửa vỏ. Một tiếng *TÁCH* khẽ vang lên ngắn gọn, gọn ghẽ, dứt khoát.

``Xong rồi\dots'' chị thì thầm, mắt mở to. ``Một chiếc đồng hồ hoàn chỉnh.''

Tôi cầm nó lên, cảm giác nặng trĩu và lạnh giá lan từ lòng bàn tay lên cánh tay. Mặt đồng hồ vẫn tối đen, không có gì hiển thị. Không ánh sáng, không tiếng tích tắc.

Nhưng vấn đề là\dots

``Không có ổ\dots\ Không có nguồn.'' anh nói, mặc dù tôi nhớ là Onii-sama đã cầm củ sạc và đế sạc

``Và ở đây\dots\ không chỗ để sạc,'' chị tôi tiếp lời.

``Chưa kể rằng ở đây cũng không chắc có điện nữa,'' tôi thở dài.

Tôi nhìn vào chiếc đồng hồ đã được lắp hoàn thiện, cảm thấy tim mình đập mạnh đến mức đau. ``Chúng ta thực sự là ai? Và cái thứ này\dots\ sẽ cho ta câu trả lời gì?''

Anh tôi cúi đầu, đôi mắt kính lấp lánh phản chiếu mặt đồng hồ tối om. ``Có lẽ\dots\ khi nó sáng lên, chúng ta sẽ biết.``

``Nhưng bây giờ\dots'' tôi ôm chiếc đồng hồ vào lồng ngực, cảm giác lạnh giá thấm qua áo, ``chúng ta chỉ có thể chờ.``

Một cơn thôi thúc mãnh liệt dấy lên trong tôi: nếu lắp ghép xong, có lẽ chúng tôi sẽ hiểu ra điều gì đó. Nhưng ngay cả khi mọi thứ đã khớp vào nhau, tôi vẫn chẳng biết vai trò thật sự của nó là gì. Chỉ một cái khớp nối thôi, sao lại được giấu kỹ đến thế, và tại sao lại nằm lẫn trong cuốn kỷ yếu?

Tôi siết chặt mảnh kim loại trong tay, tim đập thình thịch. Trong đầu tôi, hai câu hỏi vang lên, rõ rệt đến mức như muốn xé rách bầu không khí im lìm:

Chúng tôi thực sự là ai? Và cái tên ``Hikawa'' gắn liền với mình có ý nghĩa gì?

Và hơn cả, chiếc đồng hồ kia, rốt cuộc được tạo ra để làm gì?

Lúc này đây, trong khoảnh khắc ấy, tôi cảm thấy một ánh sáng rất mảnh hiện lên: một lối đi duy nhất, để tìm lại bản thân thật sự của chúng tôi. Chúng tôi không được phép quay lưng lại với sự thật đang chờ phía trước.

\end{document}